\documentclass[11pt]{article}

\usepackage[margin=1in]{geometry}
\usepackage{microtype}
\usepackage{listings}
\usepackage{xcolor}
\usepackage{hyperref}
\usepackage{booktabs}
\usepackage{enumitem}

\hypersetup{
  colorlinks=true,
  linkcolor=blue!60!black,
  citecolor=blue!60!black,
  urlcolor=blue!60!black,
}

% Modula-3 listing style
\lstdefinelanguage{Modula3}{
  morekeywords={MODULE,INTERFACE,PROCEDURE,BEGIN,END,VAR,CONST,TYPE,
    IF,THEN,ELSE,ELSIF,WHILE,DO,FOR,TO,BY,REPEAT,UNTIL,LOOP,EXIT,
    RETURN,WITH,CASE,OF,LOCK,TRY,EXCEPT,FINALLY,RAISE,RAISES,
    IMPORT,FROM,REVEAL,BRANDED,OBJECT,METHODS,OVERRIDES,
    ARRAY,RECORD,SET,REF,REFANY,ROOT,UNTRACED,
    NIL,TRUE,FALSE,AND,OR,NOT,MOD,DIV,IN,
    GENERIC,UNSAFE,VALUE,READONLY,
    INTEGER,BOOLEAN,CHAR,TEXT,REAL,LONGREAL,EXTENDED,CARDINAL,MUTEX},
  sensitive=true,
  morecomment=[s]{(*}{*)},
  morestring=[d]{"},
}

% ESC annotation style
\lstdefinelanguage{ESC}{
  morekeywords={SPEC,LOOPINV,PRAGMA,
    REQUIRES,ENSURES,MODIFIES,
    VAR,MAP,TO,SEQ,
    ABSTRACT,DEPEND,FUNC,AXIOM,INVARIANT,REP,
    ALL,EX,RES,FRESH,NUMBER,
    AND,OR,NOT,IMPLIES,IFF,
    MEMBER,INSERT,DELETE,sup,
    CONCAT,SUBARRAY,MIN,MAX,
    NIL,TRUE,FALSE,INTEGER,BOOLEAN,MUTEX},
  sensitive=true,
  morecomment=[s]{(*}{*)},
}

\lstset{
  basicstyle=\ttfamily\small,
  keywordstyle=\bfseries,
  commentstyle=\itshape\color{gray},
  breaklines=true,
  frame=single,
  framerule=0.4pt,
  xleftmargin=2em,
  framexleftmargin=1.5em,
  numbers=none,
  columns=flexible,
}

\newcommand{\escm}{ESC/Modula-3}
\newcommand{\simplify}{\textsc{Simplify}}
\newcommand{\m}[1]{\texttt{#1}}
\newcommand{\kw}[1]{\texttt{\textbf{#1}}}

\title{ESC/Modula-3 Mk.\ II}
\author{Mika Nystr\"om}
\date{March 2026}

\begin{document}
\maketitle

\tableofcontents

\bigskip
\noindent
\textbf{Note.}  \escm{} is under active development.
Phase~0 (pragma extraction and specification parsing) is complete.
Later phases (verification condition generation, proving, error
reporting) are in progress.  This manual documents the annotation
language, the command-line interface, and the overall architecture.
Sections on VCG and error output will be expanded as those phases
are completed.

%======================================================================
\section{Introduction}
\label{sec:intro}

\escm{} is an Extended Static Checker for Modula-3.  It catches common
runtime errors at compile time by generating first-order logical
formulas from annotated source code and discharging them with the
\simplify{} automatic theorem prover.

Classes of bugs detected:
\begin{itemize}[nosep]
  \item Null pointer dereferences (\m{NIL} access)
  \item Array index out of bounds
  \item Type-narrowing failures (runtime \m{NARROW} errors)
  \item Race conditions and locking-order violations
  \item Violated preconditions, postconditions, and invariants
\end{itemize}

\subsection{History}

The original ESC/Modula-3 was developed at DEC Systems Research Center
(later Compaq SRC) from approximately 1991 to 1998 by David Detlefs,
Greg Nelson, James Saxe, and Rustan Leino.  The companion tool
ESC/Java~\cite{flanagan:pldi02} followed in 1998--2002.

The system had four major components: (1)~a~front-end that parsed
Modula-3 with annotations, (2)~a~translator to guarded commands,
(3)~a~verification condition generator (VCG), and (4)~the \simplify{}
theorem prover.

When HP absorbed Compaq SRC, the front-end (${\approx}\,34{,}000$
lines of Modula-3) was never released.  The \simplify{} prover, the
axiom files encoding Modula-3 semantics, and the annotated standard
library interfaces survived in the CM3 repository.

This project reconstructs the lost front-end using the surviving
artifacts as specifications.  The prover is unchanged; the VCG is
reimplemented using a Modula-3/Scheme hybrid architecture.

%======================================================================
\section{Building and Running}
\label{sec:build}

\subsection{Prerequisites}

A working CM3 compiler installation is required.  See the CM3
\m{CLAUDE.md} for bootstrap instructions.  The \m{cm3} binary
must be on \m{PATH}.

\subsection{Building \escm{}}

The system comprises four packages, built in dependency order:

\begin{lstlisting}[language={}]
cd ESC/escm3/esclib/src  && cm3 -build && cm3 -ship
cd ../../vcgen/src       && cm3 -build && cm3 -ship
cd ../../esctool/src     && cm3 -build && cm3 -ship
\end{lstlisting}

The resulting binary is
\m{esctool/\$\{TARGET\}/escm3}, where \m{TARGET} is the CM3 target
(e.g., \m{ARM64\_DARWIN}, \m{AMD64\_LINUX}).

\subsection{Building \simplify{}}

The \simplify{} prover lives in \m{ESC/Simplify/} and is built
separately:

\begin{lstlisting}[language={}]
cd ESC/Simplify/cgi-utils/src && cm3 -build && cm3 -ship
cd ../../prover/src           && cm3 -build && cm3 -ship
cd ../../simplify/src         && cm3 -build && cm3 -ship
\end{lstlisting}

\subsection{Generating a TFile}
\label{sec:tfile}

\escm{} uses \m{m3tk} to parse Modula-3 source.  It needs to find
imported interfaces (e.g., \m{Fmt.i3}, \m{Thread.i3}).  The \m{-T}
flag points to a \emph{TFile}---a text file mapping unit names to
directories, in the format consumed by \m{M3DirFindFile.TFinder}.

A TFile lists directories with lines beginning \m{@}, followed by
the filenames found in each directory.  Example:

\begin{lstlisting}[language={}]
@ /usr/local/cm3/pkg/libm3/src/fmtlex
Fmt.i3
Fmt.m3
Lex.i3
@ /usr/local/cm3/pkg/m3core/src/thread/Common
Thread.i3
\end{lstlisting}

The test suite includes a pre-populated TFile at
\m{esctest/test/m3imptab.txt}.

A TFile can be generated from an installed CM3 tree with a shell
one-liner:

\begin{lstlisting}[language={}]
M3PKG=$(dirname $(dirname $(which cm3)))/pkg
find "$M3PKG" -name '*.i3' -o -name '*.ig' -o -name '*.m3' \
  -o -name '*.mg' | sort | \
  awk -F/ '{dir=""; for(i=1;i<NF;i++) dir=dir"/"$i;
            if(dir!=prev){print "@ "dir; prev=dir} print $NF}'
\end{lstlisting}

\subsection{Running \escm{}}

\begin{lstlisting}[language={}]
escm3 -i "Fmt" -Tm3imptab.txt
escm3 -i "Sequence" "SequenceRep" -Tm3imptab.txt
\end{lstlisting}

Current output (Phase~0) lists extracted annotations with source
locations:

\begin{lstlisting}[language={}]
ESC/Modula-3 -- Extended Static Checker
Scanning for ESC annotations...
  Fmt:
    SPEC [19:1] Bool(b) ENSURES RES # NIL
    SPEC [20:1] Int(n, base) ENSURES RES # NIL
    ...
2 unit(s) processed, 14 annotation(s) found.
\end{lstlisting}

\subsection{Running the test suite}

\begin{lstlisting}[language={}]
cd ESC/escm3/esctest/test
make test OUTPUTDIR=ARM64_DARWIN
\end{lstlisting}

Available targets: \m{test-fmt} (pragma extraction),
\m{test-list} (interface checking), \m{test-seq} (VC generation),
\m{test-golden} (golden VC comparison).

%======================================================================
\section{CLI Reference}
\label{sec:cli}

\escm{} is built on the \m{M3ToolFrame}/\m{M3CFETool} framework from
\m{m3tk}.  It inherits the standard m3tk command-line flags.

\subsection{Primary flags}

\begin{table}[ht]
\centering
\begin{tabular}{@{}lll@{}}
\toprule
Flag & Argument & Description \\
\midrule
\m{-i} & \m{"Name1" "Name2" \ldots} & Interfaces to check \\
\m{-m} & \m{"Name1" "Name2" \ldots} & Modules to check \\
\m{-T}\emph{file} & (prefix) & TFile for source lookup \\
\m{-D}\emph{path} & (prefix) & Search directory for sources \\
\m{-pn} & \m{"path1" "path2" \ldots} & Explicit source file paths \\
\m{-NOSTD} & (flag) & Suppress default library paths \\
\bottomrule
\end{tabular}
\caption{Command-line flags inherited from \m{m3tk}.}
\label{tab:flags}
\end{table}

The \m{-i} (Interfaces) and \m{-m} (Modules) flags name compilation
units by their Modula-3 unit name (not filename).  Multiple names are
given as separate quoted arguments.

The \m{-T} and \m{-D} flags are \emph{prefix} arguments: the value is
concatenated directly after the flag letter with no space.  For
example, \m{-T/tmp/m3imptab.txt} or \m{-D/usr/local/cm3/pkg/libm3/src}.
Multiple occurrences are allowed.

The \m{-pn} (PathNames) flag provides explicit paths to source files
as positional arguments.

%======================================================================
\section{Annotation Language Reference}
\label{sec:annotations}

\escm{} annotations are embedded in Modula-3 source as pragmas.
Two pragma keywords are used: \kw{SPEC} for specifications and
\kw{LOOPINV} for loop invariants.

\subsection{Enabling pragmas}

Before any annotations can appear, the pragma keyword must be
declared:

\begin{lstlisting}[language=Modula3]
<*PRAGMA SPEC*>
<*PRAGMA SPEC, LOOPINV*>
\end{lstlisting}

This is standard Modula-3 syntax: \m{PRAGMA} declares user-defined
pragma keywords so the parser will not reject them.

\subsection{Procedure specifications}
\label{sec:procspec}

A procedure specification declares preconditions, postconditions, and
frame conditions:

\begin{lstlisting}[language=ESC]
<*SPEC proc(arg1, arg2, ...)
       REQUIRES P
       ENSURES  Q
       MODIFIES V1, V2, ... *>
\end{lstlisting}

\begin{description}[nosep,leftmargin=2em]
  \item[\kw{REQUIRES}] Precondition---the caller's obligation.  Zero
    or more clauses; implicitly conjoined.
  \item[\kw{ENSURES}] Postcondition---the callee's guarantee.  Zero
    or more clauses; implicitly conjoined.
  \item[\kw{MODIFIES}] Frame condition---the variables or map entries
    the procedure may change.  Everything not listed is guaranteed
    unchanged.
\end{description}

The formal parameter list names the arguments for use in the
specification body.  For object methods, the receiver is the first
argument.

\medskip\noindent\textbf{Example.}

\begin{lstlisting}[language=ESC]
<*SPEC T.addhi(t, x)
       MODIFIES Data[t]
       REQUIRES Valid[t]
       ENSURES NUMBER(Data'[t]) = NUMBER(Data[t])+1
           AND (ALL [i: INTEGER] 0 <= i AND i < NUMBER(Data[t])
                  IMPLIES Data'[t][i] = Data[t][i])
           AND Data'[t][NUMBER(Data[t])] = x *>
\end{lstlisting}

This specifies that \m{addhi} requires the sequence to be valid,
modifies only the abstract data, and ensures the new data is the old
data with element~\m{x} appended.

\subsection{Ghost variables}
\label{sec:ghost}

Ghost variables exist only in specifications; they have no runtime
representation.

\medskip\noindent\textbf{Plain ghost variable.}

\begin{lstlisting}[language=ESC]
<*SPEC VAR Valid: BOOLEAN *>
\end{lstlisting}

\medskip\noindent\textbf{Mapped variable (per-object abstract state).}

\begin{lstlisting}[language=ESC]
<*SPEC VAR Valid: MAP T TO BOOLEAN *>
<*SPEC VAR Data: MAP T TO SEQ[Elem.T] *>
\end{lstlisting}

A \kw{MAP} variable associates each object of type~\m{T} with a value
of the range type.  \m{Valid[t]} denotes the value for object~\m{t}.
A \kw{SEQ} range type models a variable-length sequence.

\subsection{Abstraction functions}

An abstraction function defines a ghost variable in terms of concrete
fields:

\begin{lstlisting}[language=ESC]
<*SPEC ABSTRACT Seq.Valid[t: Seq.T]:
                Seq.Valid[t] IFF
                      t # NIL AND t.elem # NIL
                  AND NUMBER(t.elem^) > 0
                  AND t.st < NUMBER(t.elem^)
                  AND t.sz <= NUMBER(t.elem^) *>
\end{lstlisting}

This definition appears in a \kw{REVEAL} section (e.g.,
\m{SequenceRep.ig}) where the concrete representation is visible.
The name is fully qualified: \m{Module.Variable}.

\subsection{Abstraction dependencies}

A \kw{DEPEND} declaration lists the concrete fields that an abstract
variable depends on.  This information is used for frame-condition
reasoning: modifying a concrete field can affect only those abstract
variables that depend on it.

\begin{lstlisting}[language=ESC]
<*SPEC DEPEND Seq.Data[t: Seq.T]:
              t.st, t.sz, t.elem, RefArray *>
\end{lstlisting}

\subsection{Specification functions}

A specification-only function can be used in annotations but has no
runtime code:

\begin{lstlisting}[language=ESC]
<*SPEC FUNC m(i: INTEGER, n: INTEGER): INTEGER *>
\end{lstlisting}

Semantics are given by axioms (Section~\ref{sec:axioms}).

\subsection{Axioms}
\label{sec:axioms}

An axiom asserts a logical fact, typically to define specification
functions:

\begin{lstlisting}[language=ESC]
<*SPEC AXIOM (ALL [i: INTEGER, n: INTEGER]
                i >= n IMPLIES m(i, n) = i - n) *>
<*SPEC AXIOM (ALL [i: INTEGER, n: INTEGER]
                i < n IMPLIES m(i, n) = i) *>
\end{lstlisting}

Axioms are implicitly universally quantified at the top level and are
added to the prover's background theory.

\subsection{Object invariants}

An invariant must hold at every stable program state:

\begin{lstlisting}[language=ESC]
<*SPEC INVARIANT (ALL [t1: Seq.T, t2: Seq.T]
                    t1 # t2 IMPLIES t1.elem # t2.elem) *>
\end{lstlisting}

\subsection{Representation invariants}

A \kw{REP} declaration specifies a representation invariant for use in
\kw{REVEAL} sections:

\begin{lstlisting}[language=ESC]
<*SPEC REP P *>
\end{lstlisting}

\subsection{Loop invariants}

Loop invariants use the separate \kw{LOOPINV} pragma keyword and
appear immediately before a \m{WHILE} loop:

\begin{lstlisting}[language=Modula3]
<*LOOPINV i <= length AND sorted(a, 0, i) *>
WHILE i < length DO
  ...
END;
\end{lstlisting}

The invariant must hold on entry to the loop, be preserved by each
iteration, and is assumed to hold on loop exit (conjoined with the
negated loop condition).

\subsection{Expression sub-language}
\label{sec:exprs}

Annotations use a first-order expression language embedded in Modula-3
pragmas.  The following table lists all operators and built-in
expressions, from lowest to highest precedence.

\begin{table}[ht]
\centering
\begin{tabular}{@{}lll@{}}
\toprule
Operator / Expression & Meaning & Assoc. \\
\midrule
\m{E1 IFF E2} & Logical equivalence & --- \\
\m{E1 IMPLIES E2}, \m{E1 => E2} & Implication & right \\
\m{E1 OR E2} & Disjunction & left \\
\m{E1 AND E2} & Conjunction & left \\
\m{NOT E} & Negation & prefix \\
\m{E1 = E2}, \m{E1 \# E2} & Equality, inequality & --- \\
\m{E1 < E2}, \m{E1 <= E2}, \m{E1 > E2}, \m{E1 >= E2}
  & Relational comparison & --- \\
\m{E1 + E2}, \m{E1 - E2} & Addition, subtraction & left \\
\m{E1 * E2} & Multiplication & left \\
\m{- E} & Unary minus & prefix \\
\m{E.field} & Field selection & postfix \\
\m{E[index]} & Array/map indexing & postfix \\
\m{E\textasciicircum} & Dereference & postfix \\
\m{V'} & Post-state (primed variable) & postfix \\
\bottomrule
\end{tabular}
\caption{Operators by precedence (lowest first).}
\label{tab:ops}
\end{table}

\medskip\noindent\textbf{Quantifiers.}

\begin{lstlisting}[language=ESC]
ALL [x: T] P          (* universal *)
ALL [x: T, y: U] P    (* multi-variable *)
EX  [x: T] P          (* existential *)
\end{lstlisting}

\medskip\noindent\textbf{Built-in expressions.}

\begin{table}[ht]
\centering
\begin{tabular}{@{}ll@{}}
\toprule
Expression & Meaning \\
\midrule
\m{RES} & Return value of enclosing function \\
\m{V'[x]} or \m{V'} & Post-state value (after procedure execution) \\
\m{FRESH(x)} & \m{x} is freshly allocated \\
\m{NUMBER(a)} & Length of array or sequence \m{a} \\
\m{SUBARRAY(a, i, n)} & Subarray of \m{a} starting at \m{i}, length \m{n} \\
\m{CONCAT(a, b)} & Sequence concatenation \\
\m{MIN(a, b)}, \m{MAX(a, b)} & Minimum, maximum \\
\midrule
\multicolumn{2}{@{}l}{\emph{Set operations (used for lock levels):}} \\
\m{MEMBER(x, S)} & \m{x} is a member of set \m{S} \\
\m{INSERT(x, S)} & Set \m{S} with \m{x} added \\
\m{DELETE(x, S)} & Set \m{S} with \m{x} removed \\
\m{sup(S)} & Supremum (maximum element) of \m{S} \\
\midrule
\multicolumn{2}{@{}l}{\emph{Constants:}} \\
\m{NIL} & Null reference \\
\m{TRUE}, \m{FALSE} & Boolean constants \\
\m{0}, \m{-1}, \m{42}, \ldots & Integer literals \\
\bottomrule
\end{tabular}
\caption{Built-in expressions.}
\label{tab:builtins}
\end{table}

\medskip\noindent\textbf{Lock level variable.}\enspace The implicit ghost variable
\m{LL} denotes the set of locks currently held by a thread.  It is
used with the set operations to specify locking discipline
(Section~\ref{sec:thread}).

%======================================================================
\section{Worked Examples}
\label{sec:examples}

This section walks through four annotated library interfaces from the
CM3 standard library.  These annotations were written by the original
ESC/Modula-3 authors (Detlefs, Nelson, Heydon, McJones) and survive
in the CM3 source tree.

\subsection{Fmt.i3 --- Simple postconditions}
\label{sec:fmt}

\begin{lstlisting}[language=ESC]
<* PRAGMA SPEC *>

<*SPEC Bool(b)       ENSURES RES # NIL *>
<*SPEC Int(n, base)  ENSURES RES # NIL *>
<*SPEC Real(x, ...)  ENSURES RES # NIL *>
\end{lstlisting}

Every formatting procedure guarantees its result is non-nil.  The
annotation \m{ENSURES RES \# NIL} means: after the call returns, the
result reference is not \m{NIL}.  This lets the checker suppress
false null-dereference warnings at call sites that immediately use the
result (e.g., \m{Wr.PutText(wr, Fmt.Int(n))}).

There are no preconditions (\kw{REQUIRES}) and no frame conditions
(\kw{MODIFIES}), because formatting is a pure function with no
side effects on abstract state.

\subsection{List.ig --- Allocation and field access}
\label{sec:list}

\begin{lstlisting}[language=ESC]
<* PRAGMA SPEC *>

<*SPEC Cons(head, tail)
       ENSURES RES # NIL AND FRESH(RES)
           AND RES.head = head AND RES.tail = tail *>
\end{lstlisting}

The \m{Cons} specification says:
\begin{enumerate}[nosep]
  \item The result is not \m{NIL}.
  \item The result is \kw{FRESH}---newly allocated, distinct from any
    pre-existing object.  This is crucial for the checker to know that
    the new cons cell does not alias any existing data.
  \item The \m{head} and \m{tail} fields of the result equal the
    corresponding arguments.
\end{enumerate}

The \kw{FRESH} predicate is central to reasoning about allocation.
Without it, the checker could not rule out aliasing between the new
cell and cells already in the list.

\subsection{Sequence.ig + SequenceRep.ig --- Full data abstraction}
\label{sec:seq}

This is the richest example, demonstrating ghost variables, mapped
types, abstraction functions, specification functions, axioms, and
dependencies.

\medskip\noindent\textbf{Public interface (Sequence.ig).}

\begin{lstlisting}[language=ESC]
<* PRAGMA SPEC *>

<*SPEC VAR Valid: MAP T TO BOOLEAN *>
<*SPEC VAR Data:  MAP T TO SEQ[Elem.T] *>
\end{lstlisting}

Two ghost variables model the abstract state of a sequence object:
\m{Valid[t]} is a boolean saying whether \m{t} is a well-formed
sequence, and \m{Data[t]} is the abstract sequence of elements.

\begin{lstlisting}[language=ESC]
<*SPEC T.addhi(t, x)
       MODIFIES Data[t]
       REQUIRES Valid[t]
       ENSURES NUMBER(Data'[t]) = NUMBER(Data[t])+1
           AND (ALL [i: INTEGER] 0 <= i AND i < NUMBER(Data[t])
                  IMPLIES Data'[t][i] = Data[t][i])
           AND Data'[t][NUMBER(Data[t])] = x *>

<*SPEC T.size(t)
       REQUIRES Valid[t]
       ENSURES RES = NUMBER(Data[t]) *>

<*SPEC T.gethi(t)
       REQUIRES Valid[t] AND NUMBER(Data[t]) > 0
       ENSURES RES = Data[t][NUMBER(Data[t])-1] *>
\end{lstlisting}

Key observations:
\begin{itemize}[nosep]
  \item \m{MODIFIES Data[t]} restricts the procedure to changing only
    the data of the specific sequence object~\m{t}.  All other
    sequences, and \m{Valid[t]} itself, are guaranteed unchanged.
  \item \m{Data'[t]} denotes the post-state of \m{Data[t]}.  The
    primed notation makes pre/post reasoning explicit.
  \item \m{NUMBER(Data[t])} is the length of the abstract sequence.
  \item The \kw{ALL} quantifier expresses that all existing elements
    are preserved.
\end{itemize}

\medskip\noindent\textbf{Representation module (SequenceRep.ig).}

\noindent This interface reveals the concrete type and defines abstraction
functions linking ghost state to real fields:

\begin{lstlisting}[language=ESC]
<*SPEC ABSTRACT Seq.Valid[t: Seq.T]:
                Seq.Valid[t] IFF
                      t # NIL AND t.elem # NIL
                  AND NUMBER(t.elem^) > 0
                  AND t.st < NUMBER(t.elem^)
                  AND t.sz <= NUMBER(t.elem^) *>

<*SPEC DEPEND Seq.Data[t: Seq.T]:
              t.st, t.sz, t.elem, RefArray *>
\end{lstlisting}

The \kw{ABSTRACT} declaration says: ``\m{Valid[t]} is true if and only
if the concrete fields satisfy these constraints.''  This lets the
checker verify procedure bodies against their abstract specifications.

The \kw{DEPEND} declaration lists which concrete fields the abstract
variable \m{Data} depends on.  When a procedure modifies \m{t.st},
the checker knows only \m{Data[t]} (and other abstract variables
depending on \m{t.st}) might change.

\begin{lstlisting}[language=ESC]
<*SPEC FUNC m(i: INTEGER, n: INTEGER): INTEGER *>
<*SPEC AXIOM (ALL [i: INTEGER, n: INTEGER]
                i >= n IMPLIES m(i, n) = i - n) *>
<*SPEC AXIOM (ALL [i: INTEGER, n: INTEGER]
                i < n IMPLIES m(i, n) = i) *>

<*SPEC ABSTRACT Seq.Data[t: Seq.T]:
       Seq.Data[t] = Abs(t.elem^, t.st, t.sz) *>
\end{lstlisting}

The specification function \m{m} is a modular subtraction
helper defined axiomatically.  It exists only in the specification
world---there is no runtime code for it.

\subsection{Thread.i3 --- Locking discipline}
\label{sec:thread}

\begin{lstlisting}[language=ESC]
<* PRAGMA SPEC *>

<* SPEC FUNC MaxLL(m: MUTEX): BOOLEAN *>
<* SPEC AXIOM (ALL [m1, m2: MUTEX]
                   (NOT MaxLL(m1) AND MaxLL(m2)
                    AND m1 # NIL AND m2 # NIL)
                   IMPLIES m1 < m2) *>

<* SPEC Acquire(m) MODIFIES LL
                   REQUIRES sup(LL) < m
                   ENSURES LL' = INSERT(LL, m) *>

<* SPEC Release(m) MODIFIES LL
                   REQUIRES MEMBER(m, LL)
                   ENSURES LL' = DELETE(LL, m) *>
\end{lstlisting}

This is ESC's locking discipline, which prevents deadlocks by
enforcing a total order on lock acquisition:

\begin{itemize}[nosep]
  \item \m{LL} is the implicit ghost variable representing the set of
    locks currently held by the executing thread.
  \item \m{Acquire(m)} requires \m{sup(LL) < m}: you may only acquire
    a lock that is ``higher'' than all currently held locks.  This
    prevents circular wait.
  \item \m{Release(m)} requires \m{MEMBER(m, LL)}: you can only
    release a lock you actually hold.
  \item \m{MaxLL(m)} marks a mutex as having the maximal lock level,
    above all non-maximal mutexes.
  \item Both operations modify only \m{LL}.
\end{itemize}

This simple protocol is sufficient to prevent deadlocks in programs
that follow the locking discipline.  The checker verifies compliance
statically.

%======================================================================
\section{Architecture Overview}
\label{sec:arch}

\subsection{Package structure}

\begin{lstlisting}[language={}]
ESC/escm3/
  esclib/src/         Core library (7 modules)
    ESCPragma.m3        Pragma extraction from m3tk AST
    ESCSpec.i3          AST types for parsed specifications
    ESCSpecParse.m3     Recursive-descent spec parser (~980 lines)
    ESCTypes.m3         Type environment from m3tk semantics
    ESCNameMap.m3       M3-to-Simplify naming conventions
    ESCSx.m3            S-expression builder for formulas
    ESCError.m3         Error reporting types

  vcgen/src/          Scheme-based VC generator (9 modules)
    esc-util.scm        S-expression constructors
    esc-gcmd.scm        Guarded command translation
    esc-passify.scm     Assignment passification
    esc-wp.scm          Weakest-precondition calculus
    esc-typegen.scm     Type predicate/axiom generation
    esc-framegen.scm    Frame condition axioms
    esc-absfunc.scm     Abstraction function constraints
    esc-bg.scm          Background predicate management
    esc-emit.scm        Labeled formula serialization

  esctool/src/        CLI driver (4 modules)
    ESCMain.m3          Entry point (M3ToolFrame-based)
    ESCDriver.m3        Pipeline orchestrator
    ESCProve.m3         Simplify prover wrapper
    ESCReport.m3        Counterexample-to-source mapping

  esctest/test/       Test suite
    Makefile            Test targets
    m3imptab.txt        TFile for package resolution

ESC/Simplify/         Simplify theorem prover (75 modules)
  prover/src/           Core prover library
  simplify/src/         CLI driver and test suites
  cgi-utils/src/        CGI/HTML utilities
\end{lstlisting}

\subsection{Data flow}

The verification pipeline has five stages:

\begin{lstlisting}[language={}]
  M3 source + <*SPEC*> annotations
         |
    [1. ESCPragma]     Extract raw pragma text from m3tk AST
         |
    [2. ESCSpecParse]  Parse spec sub-language -> Spec AST
         |
    [3. esc-gcmd]      Translate M3 body + specs -> guarded commands
         |
    [4. esc-wp]        Weakest precondition -> first-order formula
         |
    [5. Simplify]      Prove/refute -> Valid / Invalid + labels
         |
    [6. ESCReport]     Map labels -> source-level warnings
\end{lstlisting}

Stages 1--2 are implemented in Modula-3 (\m{esclib}).  Stages 3--4
are implemented in Scheme (\m{vcgen}), leveraging mscheme's zero-copy
interop with Modula-3 data structures.  Stage~5 uses the existing
\simplify{} prover.  Stage~6 maps counterexample labels back to
source locations.

\subsection{How \simplify{} fits in}

\simplify{} is a Nelson-Oppen combination of decision procedures:
\begin{itemize}[nosep]
  \item \textbf{Congruence closure} for uninterpreted functions and
    equality
  \item \textbf{Simplex} for linear arithmetic over integers and
    rationals
  \item \textbf{Partial orders} for lock-level reasoning
  \item \textbf{Pattern-driven quantifier instantiation} using
    user-supplied \m{PATS}/\m{NOPATS} hints
\end{itemize}

Input is an S-expression formula.  Output is \m{Valid} (the formula is
a tautology modulo the background theory), \m{Invalid} (with a
counterexample context and labels identifying which assertion failed),
or \m{Unknown} (resource limits exceeded).

The VCG wraps each assertion with a label of the form
\m{|ERROR.kind.line.col|}.  When \simplify{} reports \m{Invalid}, the
labels in the counterexample identify the source location and error
kind (e.g., null dereference, array bounds, precondition violation).

\subsection{Naming conventions}

The \m{ESCNameMap} module translates Modula-3 identifiers to
\simplify{}'s pipe-delimited naming convention:

\begin{table}[ht]
\centering
\begin{tabular}{@{}ll@{}}
\toprule
Modula-3 concept & \simplify{} name \\
\midrule
Field \m{T.f} in module \m{M} & \m{|M.T.f|} \\
Type predicate for \m{T} & \m{|Is\$T|} \\
Spec function \m{f} in \m{M} & \m{|FUNC.M.f|} \\
Type code for \m{T} & \m{|T.TYPECODE|} \\
Null & \m{|\$NIL|} \\
Boolean true/false & \m{|@true|}, \m{|@false|} \\
\bottomrule
\end{tabular}
\caption{Naming conventions.}
\end{table}

\subsection{The M3/Scheme hybrid}

The VCG uses a hybrid architecture:
\begin{itemize}[nosep]
  \item \textbf{Modula-3} handles: parsing (m3tk), type environment
    construction, prover invocation, error reporting.
  \item \textbf{Scheme} (via mscheme) handles: symbolic formula
    construction, guarded command translation, weakest-precondition
    calculus, S-expression serialization.
\end{itemize}

This division plays to each language's strengths.  Scheme's
quasiquoting and pattern matching are natural for building and
transforming logical formulas.  Modula-3's strong typing and
efficient data structures suit the compiler infrastructure and prover.

The mscheme interop is zero-copy: Scheme values are Modula-3
\m{REFANY} values, so m3tk AST nodes pass through Scheme unchanged
without marshaling.

%======================================================================
\section{Current Status and Roadmap}
\label{sec:roadmap}

\subsection{Completed}

\begin{description}[nosep,leftmargin=2em]
  \item[Phase 0: Pragma extraction.]
    \m{ESCPragma} walks the m3tk AST and extracts \kw{SPEC} and
    \kw{LOOPINV} pragmas with source locations.  Working and tested.

  \item[Phase 1: Specification parsing.]
    \m{ESCSpecParse} implements a recursive-descent parser for the
    full annotation language (Section~\ref{sec:annotations}).
    Handles procedure specs, ghost variables, abstraction functions,
    dependencies, spec functions, axioms, invariants, and loop
    invariants.  Working and tested against annotated library
    interfaces.
\end{description}

\subsection{In progress}

\begin{description}[nosep,leftmargin=2em]
  \item[Phase 2: Guarded command translation.]
    Translate Modula-3 procedure bodies and their specifications into
    guarded commands.  Scheme module \m{esc-gcmd.scm}.

  \item[Phase 3: VC generation.]
    Apply Flanagan-Saxe compact weakest-precondition calculus to
    produce near-linear verification conditions.  Scheme module
    \m{esc-wp.scm}.

  \item[Phase 4: Prover integration.]
    Wire up \m{ESCProve} to invoke \simplify{} on generated VCs,
    parse results, and pass counterexample labels to the reporter.

  \item[Phase 5: Error reporting.]
    Map counterexample labels to human-readable warnings with source
    locations.  \m{ESCReport} and \m{esc-emit.scm}.

  \item[Phase 6: End-to-end verification.]
    Verify annotated CM3 library modules (Sequence, List, Thread)
    against their specifications.
\end{description}

%======================================================================
\section{Environment Variables}
\label{sec:env}

The \simplify{} prover reads environment variables at startup to
control resource limits and search heuristics.

\begin{table}[ht]
\centering
\begin{tabular}{@{}llp{7cm}@{}}
\toprule
Variable & Default & Effect \\
\midrule
\m{PROVER\_KILL\_TIME} & $10^9$ & Max seconds per conjecture \\
\m{PROVER\_KILL\_ITER} & MAXINT & Max search iterations per conjecture \\
\m{PROVER\_MATCH\_DEPTH} & 3 & Max non-unit matching depth \\
\m{PROVER\_MAX\_FNUR} & 100 & Max restricted non-unit matching rounds \\
\m{PROVER\_CC\_LIMIT} & 1 & Max counterexample contexts shown \\
\m{PROVER\_NO\_PLUNGE} & (unset) & Disable depth-1 plunging \\
\m{PROVER\_LIT\_SPLIT} & (unset) & Enable literal splitting heuristic \\
\bottomrule
\end{tabular}
\caption{Key \simplify{} environment variables.
See \m{Prover.i3} for the full list (${\approx}\,40$ variables).}
\end{table}

%======================================================================
\begin{thebibliography}{99}

\bibitem{detlefs:src159}
D.~L.~Detlefs, K.~R.~M.~Leino, G.~Nelson, J.~B.~Saxe.
\newblock Extended Static Checking.
\newblock SRC Research Report 159, Compaq, 1998.

\bibitem{leino:cc98}
K.~R.~M.~Leino, G.~Nelson.
\newblock An Extended Static Checker for Modula-3.
\newblock In \emph{Compiler Construction (CC'98)}, LNCS 1383, 1998.

\bibitem{flanagan:pldi02}
C.~Flanagan, K.~R.~M.~Leino, M.~Lillibridge, G.~Nelson, J.~B.~Saxe,
R.~Stata.
\newblock Extended Static Checking for Java.
\newblock In \emph{PLDI}, 2002.

\bibitem{flanagan:popl01}
C.~Flanagan, J.~B.~Saxe.
\newblock Avoiding Exponential Explosion: Generating Compact
Verification Conditions.
\newblock In \emph{POPL}, 2001.

\bibitem{detlefs:jacm05}
D.~L.~Detlefs, G.~Nelson, J.~B.~Saxe.
\newblock Simplify: A Theorem Prover for Program Checking.
\newblock \emph{JACM}, 52(3), 2005.

\bibitem{leino:oopsla98}
K.~R.~M.~Leino.
\newblock Data Groups: Specifying the Modification of Extended State.
\newblock In \emph{OOPSLA}, 1998.

\bibitem{leino:toplas02}
K.~R.~M.~Leino, G.~Nelson.
\newblock Data Abstraction and Information Hiding.
\newblock \emph{TOPLAS}, 24(5), 2002.

\end{thebibliography}

\end{document}
